\begin{cabstract}
    随着互联网技术的飞速发展，网络流量呈现出海量化、复杂化及高动态性的特征，
    网络安全威胁日益严峻。传统的网络异常流量检测方法往往依赖于完备的标签数据
    ，难以应对实际网络环境中普遍存在的流量“长尾分布”特性以及日益增多的“未知
    攻击”威胁。针对现有方法在稀疏正常样本覆盖不足及复杂分布拟合能力受限等问题，
    本文以无监督学习为核心技术路线，深入研究了基于加权超球体学习和多模式聚类的异
    常检测算法，并设计实现了一套网络流量异常数据检测系统。本文的主要研究内容如下：

    首先，针对正常业务流量中的长尾分布导致传统单分类模型特征中心
    偏移及高误报率问题，提出了一种基于实例多样性加权的无监督超球体学习方法。该方法引
    入实例欧氏距离作为无监督的多样性度量指标，设计了
    多样性感知加权损失函数，通过动态调节不同密度样本对模型训练的贡献度，引导超球体
    边界自适应扩张。同时，引入指数滑动平均策略平滑更新特征中心，有效缓解了由
    高频背景流量主导的中心偏移问题，显著提升了模型对长尾分布下稀疏正常流量的包容性与
    检测精度。

    其次，针对网络流量中正常行为模式多样且异常类型未知的特点，提出了一种面向未知攻
    击的多模式无监督聚类检测框架。该框架打破了传统方法采用单一聚类中心的建模局限，
    通过构建多个潜在的正常行为模式中心，共同描述网络流量的复杂分布结构。模型基于样
    本到各中心的最短距离构建判别机制，实现了对多模态正常流量的精细化拟合。实验结果
    表明，该框架在无需先验异常标签的情况下，能够有效识别偏离正常模式的未知攻击，增
    强了模型在复杂网络环境下的泛化能力与鲁棒性。

    最后，为了验证算法的工程可行性，设计并实现了一套集数据实时采集、自动化预处理、
    模型在线推理及可视化告警于一体的网络流量异常数据检测系统。系统基于FastAPI框
    架、Kafka消息队列及ElasticSearch搜索引擎等现代技术栈构建。经Apache JMeter
    高并发性能测试，系统在模拟千级并发用户的高负载场景下，依然能够保持毫秒级的响应延
    迟与稳定的吞吐量，证明了本文所提算法及系统架构在实际网络安全运维中的应用价值。
    
    \ckeywords{网络流量异常检测；无监督学习；长尾分布；多聚类中心}
\end{cabstract}

\begin{eabstract}
    With the rapid development of Internet technology, network traffic exhibits characteristics of massive volume, complexity, and high dynamism,
    The threat of cybersecurity is becoming increasingly severe. Traditional methods for detecting abnormal network traffic often rely on comprehensive labeled data
    It is difficult to cope with the ubiquitous "long-tail distribution" characteristic of traffic in actual network environments, as well as the increasing number of "unknowns"
    "Attack" threat. Addressing the issues of insufficient coverage of sparse normal samples and limited ability to fit complex distributions in existing methods,
    This article takes unsupervised learning as the core technical approach, and conducts in-depth research on the differences based on weighted hypersphere learning and multi-mode clustering
    We frequently evaluate algorithms and have designed and implemented a network traffic anomaly detection system. The main research contents of this paper are as follows:

    Firstly, the long-tail distribution in normal business traffic leads to the feature center of traditional single-classification models
    To address the issues of bias and high false alarm rate, an unsupervised hypersphere learning method based on instance diversity weighting is proposed. This method introduces
    The Euclidean distance of instances is designed as an unsupervised diversity metric
    The diversity-aware weighted loss function dynamically adjusts the contribution of samples with different densities to model training, guiding the hypersphere
    Boundary adaptive expansion. At the same time, the introduction of an exponential moving average strategy to smoothly update the feature center effectively mitigates the impact of
    The center-biased problem dominated by high-frequency background traffic significantly enhances the model's inclusiveness and robustness towards sparse normal traffic under long-tail distribution
    Detection accuracy.

    Secondly, given the diverse normal behavior patterns and unknown types of anomalies in network traffic, a method oriented towards unknown attacks is proposed
    A multi-mode unsupervised clustering detection framework. This framework breaks the modeling limitations of traditional methods that adopt a single clustering center,
    By constructing multiple potential centers of normal behavior patterns, the complex distribution structure of network traffic is jointly described. The model is based on samples
    A mechanism for determining the shortest distance to each center has been established, achieving a refined fitting of multimodal normal traffic. experimental results
    It demonstrates that the framework can effectively identify unknown attacks deviating from the normal pattern without requiring prior anomaly labels, and enhance
    It enhances the generalization ability and robustness of the model in complex network environments.

    Finally, to verify the engineering feasibility of the algorithm, a set of functions including real-time data collection, automated preprocessing
    A network traffic anomaly detection system integrating online model inference and visual alerting. The system is built on the FastAPI framework
    Built with modern technology stacks such as Docker, Kafka message queue, and ElasticSearch search engine. Tested with Apache JMeter
    High concurrency performance testing: The system can still maintain millisecond-level response latency under high load scenarios simulating thousands of concurrent users
    The late and stable throughput demonstrates the practical application value of the algorithm and system architecture proposed in this paper in network security operation and maintenance.

    \ckeywords{Network traffic anomaly detection; Unsupervised learning; Long-tail distribution; Multiple clustering centers}
\end{eabstract}
